\chapter{Технологическая часть}

\section{Средства реализации}

Компилятор реализован на языке C. Выбор C обусловлен близостью к системному программированию, удобством работы с файлами, указателями и структурами данных, а также соответствием тематике компилятора исторического языка B \cite{kandr}.

Целевая платформа:

\begin{itemize}
  \item операционная система --- Linux;
  \item архитектура --- x86-64;
  \item формат промежуточного кода --- assembly GNU assembler \cite{gas};
  \item сборка компилятора --- \texttt{make};
  \item сборка результата --- \texttt{as} и \texttt{cc} \cite{gcc}.
\end{itemize}

Командная строка компилятора поддерживает следующие режимы:

\begin{itemize}
  \item \texttt{-S} --- сгенерировать ассемблерный файл и остановиться;
  \item \texttt{-c} --- сгенерировать объектный файл и не выполнять компоновку;
  \item \texttt{-o file} --- задать имя исполняемого файла;
  \item \texttt{-T} --- вывести абстрактное синтаксическое дерево;
  \item \texttt{-v} --- вывести команды этапов компиляции.
\end{itemize}

\section{Особенности реализации лексического анализатора}

Лексический анализатор реализован в файле \texttt{scan.c}. Он посимвольно читает входной файл и формирует структуру токена. Для каждого токена сохраняется его тип, строковое описание и целочисленное значение, если токен является литералом или составным присваиванием.

Числовые литералы обрабатываются функцией, которая выбирает основание системы счисления по первому символу. Если литерал начинается с \texttt{0}, значение интерпретируется как восьмеричное. Это соответствует историческому поведению B.

Строковые и символьные литералы обрабатываются с учетом escape-последовательностей языка B. Например, строка \texttt{"\%d*n"} после лексического анализа содержит символ перевода строки, а не два символа \texttt{*} и \texttt{n}.

\section{Особенности реализации синтаксического анализатора}

Разбор программы начинается с функции обработки глобальных объявлений. Для каждого внешнего определения компилятор определяет его форму по токену, следующему за именем:

\begin{itemize}
  \item \texttt{(} означает объявление или определение функции;
  \item \texttt{[} означает определение вектора;
  \item литерал или имя после глобального объекта означает список инициализаторов;
  \item \texttt{;} завершает простое определение.
\end{itemize}

Операторы разбираются в файле \texttt{stmt.c}. Для составных блоков формируется последовательность AST-узлов, связанных операцией \texttt{A\_GLUE}. Это позволяет генератору кода выполнить операторы в исходном порядке.

Выражения разбираются в файле \texttt{expr.c}. Основой является функция разбора бинарных выражений, получающая минимальный допустимый приоритет. Такой способ позволяет корректно разбирать выражения вида:

\begin{verbatim}
x = a + b * c == d ? f(1) : g(2);
\end{verbatim}

Составные присваивания языка B используют форму \texttt{x =+ 1}. Лексер распознает такую запись как присваивание с дополнительной операцией, а парсер преобразует ее в AST, эквивалентный \texttt{x = x + 1}.

\section{Особенности генерации кода}

Генератор кода создает текст ассемблерной программы. Для каждой функции формируется пролог:

\begin{verbatim}
pushq   %rbp
movq    %rsp, %rbp
\end{verbatim}

После этого стековый указатель сдвигается на размер области локальных переменных. При выходе из функции выполняется восстановление кадра и инструкция \texttt{ret}.

Для локальных переменных используются отрицательные смещения относительно \texttt{\%rbp}. Для глобальных переменных применяются RIP-relative обращения, например \texttt{name(\%rip)}. Такой код подходит для современных x86-64 Linux систем.

Вызовы внешних функций выполняются через PLT. Передача аргументов организована в соответствии с соглашением System V AMD64 ABI \cite{sysvabi}. Благодаря этому программы на B могут вызывать внешние системные функции, например \texttt{printf}, если имя объявлено через \texttt{extrn}.

\section{Пример программы на языке B}

Листинг 3.1 --- Программа на языке B, демонстрирующая автоматический вектор, указатель и индексирование.
\begin{verbatim}
main() {
    auto v 4, p;
    extrn printf;

    v[0] = 3;
    v[1] = 4;
    v[2] = 5;
    v[3] = 6;

    p = &v[1];
    printf("%d*n", *p);

    *p = 20;
    printf("%d*n", v[1]);
    printf("%d*n", 2[v]);
    printf("%d*n", *(v + 3));

    return 0;
}
\end{verbatim}

В программе создается автоматический вектор \texttt{v} из четырех машинных слов и переменная \texttt{p}. В переменную \texttt{p} записывается адрес элемента \texttt{v[1]}. Далее через разыменование изменяется значение этого элемента. Выражения \texttt{v[2]}, \texttt{2[v]} и \texttt{*(v + 3)} демонстрируют адресную модель векторов языка B.

\section{Фрагмент полученного ассемблерного кода}

Листинг 3.2 --- Фрагмент ассемблерного кода, полученного для программы из листинга 3.1.
\begin{verbatim}
main:
    pushq   %rbp
    movq    %rsp, %rbp
    addq    $-48,%rsp
    movq    $3, %r10
    leaq    -32(%rbp), %r11
    movq    $0, %r12
    salq    $3, %r12
    addq    %r11, %r12
    movq    %r10, (%r12)
    leaq    -32(%rbp), %r10
    movq    $1, %r11
    salq    $3, %r11
    addq    %r10, %r11
    movq    %r11, -40(%rbp)
    movq    -40(%rbp), %r10
    movq    (%r10), %r10
    movq    %r10, %rsi
    leaq    L2(%rip), %r10
    movq    %r10, %rdi
    call    printf@PLT
\end{verbatim}

В приведенном фрагменте видно размещение локального вектора в стековом кадре, масштабирование индекса на 8 байт, вычисление адреса элемента и вызов внешней функции \texttt{printf}.

\section{Результат выполнения программы}

Листинг 3.3 --- Вывод программы из листинга 3.1.
\begin{verbatim}
4
20
5
6
\end{verbatim}

Первое число является исходным значением \texttt{v[1]}. Второе число показывает изменение элемента через указатель \texttt{p}. Последние две строки демонстрируют корректное чтение элементов вектора через эквивалентные формы индексирования.

\section{Пример рекурсивной программы}

Листинг 3.4 --- Рекурсивное вычисление числа Фибоначчи.
\begin{verbatim}
fib(n) {
    if (n <= 1) return(n);
    return(fib(n - 1) + fib(n - 2));
}

main() {
    extrn printf;

    printf("%d*n", fib(10));
    return(0);
}
\end{verbatim}

Листинг 3.5 --- Вывод рекурсивной программы.
\begin{verbatim}
55
\end{verbatim}

Этот пример проверяет определение функции, передачу параметров, рекурсивный вызов, условный оператор и возврат значения.

\section{Тестирование}

Для проверки компилятора подготовлен набор тестовых программ в каталоге \texttt{tests}. Каждая программа \texttt{inputXX.b} имеет соответствующий файл ожидаемого вывода \texttt{out.inputXX.b}. Скрипт \texttt{tests/runtests} компилирует каждую программу, запускает полученный исполняемый файл и сравнивает фактический вывод с эталонным.

Основные проверяемые возможности приведены в таблице \ref{tbl:tests}.

\begin{longtable}{|p{4cm}|p{12cm}|}
  \caption{Проверяемые возможности компилятора}
  \label{tbl:tests} \\
  \hline
  \textbf{Тесты} & \textbf{Проверяемые конструкции} \\
  \hline
  \endfirsthead
  \hline
  \textbf{Тесты} & \textbf{Проверяемые конструкции} \\
  \hline
  \endhead
  \texttt{input01.b} & Рекурсия, условный оператор, цикл, арифметика, вызовы функций. \\
  \hline
  \texttt{input02.b} & Сортировка массивов, вложенные циклы, сравнения, глобальные векторы. \\
  \hline
  \texttt{input03.b}, \texttt{input04.b} & Обход графа в ширину и глубину, глобальные массивы, вложенные вызовы функций. \\
  \hline
  \texttt{input05.b} & Моделирование связного списка через массивы индексов. \\
  \hline
  \texttt{input06.b}--\texttt{input09.b} & Пустые операторы, составные присваивания, отрицательные числа, инициализация глобальных объектов. \\
  \hline
  \texttt{input10.b}--\texttt{input16.b} & Указатели, адреса элементов вектора, тернарный оператор, \texttt{goto}, \texttt{switch}, вызовы функций через переменную. \\
  \hline
  \texttt{input17.b}, \texttt{input18.b} & Внешние объявления, глобальные инициализаторы, комментарии и escape-последовательности B. \\
  \hline
\end{longtable}

Команда запуска тестов:

\begin{verbatim}
make test
\end{verbatim}

Пример строки успешного прохождения теста:

\begin{verbatim}
input10.b: OK
\end{verbatim}

\section{Вывод}

В технологической части описаны средства реализации, особенности лексического и синтаксического анализа, генерации x86-64 кода, режимы командной строки, примеры программ на языке B, фрагмент полученного ассемблерного кода и подход к регрессионному тестированию.
